\section{Discussion}

\subsection{In-distribution trust versus out-of-distribution extrapolation}

The clearest message of the results is a sharp boundary between two regimes. \emph{In-distribution},
the loop is trustworthy on every axis we can test: the solver reproduces dataset pressure at
$\rho \approx 0.96$ across the whole $(Re, H)$ grid, and generated shapes hit their target morphology
at $\rho \approx 0.95$. This is a genuine, independent confirmation of inverse design --- the grading
physics shares no component with the generator, so a high correlation cannot be the model marking its
own work. \emph{Out-of-distribution}, fidelity degrades smoothly and predictably with distance from
the training manifold: morphology correlation softens from $\approx 0.96$ to $\approx 0.86$ along the
controlled $+\mathrm{PC1}$ ladder. The model interpolates within the laminar elliptical regime it has
seen and extrapolates poorly beyond it --- exactly the behaviour the independent loop was built to
expose, and a useful one to have quantified rather than assumed.

\subsection{Magnitude creep is a model-side calibration issue}

Two independent lines of evidence localise the amplitude bias to the generator, not the solver. First,
the resolution-, node-density- and roughness-controlled ladder leaves the affine slope $a$ unmoved by
any CFD or discretisation lever, with $a$ already elevated \emph{in-distribution}. Second, the
systematic OOD ladder measures $a \approx 2.0$--2.5 on generated shapes against the solver's validated
$a \approx 1.74$ on real dataset shapes run through the same pipeline --- the same solver, the same
mesh family, the same normalisation, differing only in whether the shape was generated or real. The
gap is therefore generative: the model captures pressure \emph{shape} faithfully but systematically
overshoots its \emph{amplitude}. This is why we report correlation, not absolute amplitude, as the
inverse-design verdict, and treat the slope as a calibration target to be closed model-side (for
example via the deferred physical-$Re$ conditioning) rather than a solver defect to be chased in the
CFD.

\subsection{Why posterior guidance is the wrong knob out-of-distribution}

The intuitive fix for the OOD gap --- adding a likelihood gradient from the model's own pressure head
via Diffusion Posterior Sampling \cite{chung2023dps} --- makes every metric worse, monotonically with
guidance strength. The mechanism is that DPS steers toward the \emph{self-trained} head's estimate of
the target, and out-of-distribution that head is precisely the component that is unreliable; stronger
guidance just optimises the surrogate's own error and inflates the body. This is a clean negative
result: plain classifier-free guidance is the correct OOD configuration, and the OOD gap cannot be
closed cheaply with a self-graded posterior. It must be closed by widening what the model genuinely
knows --- more diverse training data, or physically-grounded conditioning --- not by trusting the
surrogate harder.

\subsection{Limitations}

The findings above are bounded by several limitations, which also frame the next phase of work.

\paragraph{Validation framework still being verified.} All OpenLB-derived numbers are provisional
pending re-verification of the validation framework, and are flagged as such throughout. The
qualitative conclusions --- in-distribution trust, OOD decay, model-side magnitude creep, the DPS
negative --- are robust to modest shifts in the exact figures, but the precise correlations and slopes
should be read as current best estimates rather than settled values.

\paragraph{Narrow physical regime.} The entire study lives in 2D incompressible laminar flow around
elliptical cylinders at $Re$ 500--1000 (the vortex-shedding regime), with a single-family geometry
parameterised mainly by minor axis $b$. Nothing here speaks to turbulent, compressible, 3D, or
aerofoil-like flows; the demonstrated inverse-design capability is for this controlled testbed, chosen
deliberately as a tractable first target before scaling.

\paragraph{Steady-mean pressure target.} The 101 unsteady pressure timesteps per simulation are
collapsed to a single steady-mean $C_p$, discarding all temporal structure of the shedding. The model
neither conditions on nor predicts unsteady loads, so phenomena that matter aerodynamically ---
shedding frequency, fluctuating lift, phase --- are outside its scope by construction.

\paragraph{Single solver, single mesh family, no formal grid-convergence study.} Validation rests on
one lattice-Boltzmann solver with Reynolds-scaled but not formally converged resolution
($\mathrm{RES} \approx 0.2\,Re$, set for stability rather than from a grid-independence study), and
correlation is observed to taper slightly at the high-$Re$ end through genuine under-resolution. An
external benchmark (the Hron--Turek case) and a grid-convergence study remain outstanding and are
required before the solver's absolute accuracy --- as opposed to its internal consistency --- can be
claimed.

\paragraph{Conditioning is pressure-only and band-limited.} The target is encoded as just the first
$K = 8$ DCT-II modes of $C_p(x/c)$, which capture smooth low-frequency morphology but truncate sharp
pressure features; and the model conditions on $C_p$ alone, with no physical $Re$/$H$ conditioning.
Reynolds number is fixed (754) across the OOD ladder, so the present results cannot separate
$C_p$-morphology extrapolation from Reynolds extrapolation. Adding $Re$ conditioning is the deferred
next step and the natural lever on the magnitude creep.

\paragraph{OOD axis not isometric to conditioning space.} The PCA-$\sigma$ ladder only reaches genuine
OOD along $+\mathrm{PC1}$; the other principal directions stay near the training manifold regardless of
amplitude, because raw-curve PCA $\sigma$-units do not map linearly onto the DCT-conditioning space the
model sees. The controlled OOD claim is therefore strongest along $+\mathrm{PC1}$; a fully monotone,
isotropic OOD sweep requires re-parameterising the ladder directly in conditioning space.

\paragraph{Unresolved amplitude calibration.} The magnitude creep ($a$ elevated above the validated
1.74) is diagnosed but not yet fixed, so the model is an inverse designer of pressure \emph{shape}, not
of absolute pressure level. Until calibration is closed, downstream use must rely on correlation rather
than predicted magnitudes.

\paragraph{Surface roughness is cosmetic here but unquantified at scale.} Raw-versus-smoothed shapes
differ by only $\sim 0.01$ in correlation and $\sim 5\%$ in $C_d$ in the present regime, so the
characteristic boundary wobble does not corrupt the physics measured here. Whether that remains true at
higher $Re$, where thinner boundary layers are more sensitive to surface detail, is untested.

\subsection{Toward a rigorous evaluation}
\label{sec:toward-rigour}

The preliminary studies of Sections~\ref{sec:surface-quality} and earlier fixed design choices but are
not, on their own, rigorous; this section states how each would be put on a firm footing, alongside the
validation work proper.

\paragraph{Representation and connectivity.} The $k = 6$ connectivity choice rests on a single-seed
sweep on the \emph{circle} benchmark. A rigorous version would (i) repeat across several seeds and
report confidence intervals rather than point estimates, (ii) run the ablation on the \emph{aerodynamic}
dataset, where the relevant metric is the solver-graded $C_p$ correlation rather than a
radius-distribution KS distance, and (iii) sweep connectivity jointly with the layer count to rule out
confounding.

\paragraph{Surface quality.} The second-difference roughness of $r(\theta)$ is a convenient proxy, not a
standard measure. A rigorous treatment would adopt a resolution-independent curvature/spectral roughness
metric, tie it to a mesh-convergence study, and --- most importantly --- test directly whether smoothing
changes the solver-measured $C_p$/$C_d$ at the upper end of the Reynolds range, where boundary layers
are thinnest and most sensitive to surface detail. The present evidence that the wobble is cosmetic
holds only in the validated regime.

\paragraph{Out-of-distribution axis.} The PCA-$\sigma$ ladder reaches genuine OOD only along
$+\mathrm{PC1}$. A rigorous OOD protocol would define the ladder directly in the DCT-conditioning space
the model sees, so distance is monotone and isotropic in every direction, and would sample several seeds
per target to separate generative variance from systematic extrapolation error.

\paragraph{Validation framework and benchmarking.} Throughout, the OpenLB numbers are provisional
pending re-verification of the validation framework, and the solver itself has been checked only for
internal consistency (it reproduces the dataset normalisation). Establishing \emph{absolute} accuracy
requires a formal grid-convergence study and an external benchmark --- the Hron--Turek case --- which is
the immediate next validation milestone. The qualitative flow fields would likewise be promoted to
quantitative wake diagnostics (Strouhal number, separation point) compared against reference values.
